% ============================================================================
% SPANISCH LANGENTWURF - DS 2: FORTFÜHRUNG + ÜBERGANG
% ============================================================================
% Beobachtungssequenz Februar 2026
% Jahrgangsübergreifend Spätbeginner (Kl. 10, 11, 12)
% Lehrwerk: A_tope.com (Cornelsen)
% ============================================================================

\documentclass[11pt,a4paper]{article}

% ============================================================================
% PACKAGES
% ============================================================================
\usepackage[utf8]{inputenc}
\usepackage[T1]{fontenc}
\usepackage[ngerman]{babel}
\usepackage{geometry}
\usepackage{fancyhdr}
\usepackage{lastpage}
\usepackage{xcolor}
\usepackage{tcolorbox}
\usepackage{enumitem}
\usepackage{tabularx}
\usepackage{longtable}
\usepackage{array}
\usepackage{booktabs}
\usepackage{hyperref}
\usepackage{ifthen}
\usepackage{amssymb}

% ============================================================================
% KONFIGURATION
% ============================================================================

% --- TIER ---
\newcommand{\plantier}{1}

% --- FACH ---
\newcommand{\fachid}{spanisch}
\newcommand{\fach}{Spanisch}
\newcommand{\fachfarbe}{c41e3a}

% ============================================================================
% VARIABLEN
% ============================================================================

% --- Metadaten ---
\newcommand{\jahrgangsstufe}{10/11/12}
\newcommand{\schulart}{Gesamtschule}
\newcommand{\stundenthema}{Fortführung + Übergang}
\newcommand{\reihenthema}{Beobachtungssequenz -- Spätbeginner}
\newcommand{\stundennummer}{2}
\newcommand{\reihenstunden}{4}
\newcommand{\unterrichtsdatum}{06.01.2026}
\newcommand{\unterrichtszeit}{Doppelstunde}
\newcommand{\raum}{---}

% --- Lehrkraft ---
\newcommand{\lehrkraft}{N.N.}
\newcommand{\ausbildungsstand}{Lehrkraft}

% --- Lerngruppe ---
\newcommand{\lerngruppengroesse}{10}
\newcommand{\maennlich}{1}
\newcommand{\weiblich}{9}
\newcommand{\divers}{0}

% --- Kompetenzen/Niveau ---
\newcommand{\gerniveau}{A2/B1}
\newcommand{\lehrwerk}{A\_tope.com}
\newcommand{\lehrwerkseite}{S. 31--32 (A) / S. 104 (B)}

% ============================================================================
% LAYOUT
% ============================================================================
\geometry{left=2.5cm, right=2.5cm, top=2.5cm, bottom=2.5cm}
\definecolor{fach}{HTML}{\fachfarbe}
\definecolor{gray}{HTML}{666666}
\definecolor{lightgray}{HTML}{f5f5f5}
\definecolor{successgreen}{HTML}{2a7a4b}
\definecolor{warningorange}{HTML}{b35c00}

\tcbuselibrary{skins,breakable}

% Farbige Boxen
\newtcolorbox{infobox}[1][]{
    colback=lightgray,
    colframe=fach,
    fonttitle=\bfseries,
    title=#1,
    breakable
}

\newtcolorbox{scorebox}{
    colback=white,
    colframe=fach,
    fonttitle=\bfseries,
    title=Didaktik-Score,
    breakable
}

% Kopf-/Fußzeile
\pagestyle{fancy}
\fancyhf{}
\fancyhead[L]{\small\textcolor{gray}{\fach{} -- Jg. \jahrgangsstufe{} -- \stundenthema}}
\fancyhead[R]{\small\textcolor{gray}{Langentwurf Tier 1}}
\fancyfoot[C]{\small\thepage{} / \pageref{LastPage}}
\renewcommand{\headrulewidth}{0.4pt}
\renewcommand{\footrulewidth}{0pt}

% ============================================================================
% DOKUMENT
% ============================================================================
\begin{document}

% ============================================================================
% DECKBLATT
% ============================================================================
\begin{titlepage}
    \centering
    \vspace*{2cm}

    {\large\textcolor{gray}{\schulart}}\\[2cm]

    {\Huge\textcolor{fach}{\textbf{Langentwurf}}}\\[0.5cm]
    {\large Tier 1 -- Vollständig}\\[2cm]

    {\LARGE\textbf{\stundenthema}}\\[0.5cm]
    {\large im Rahmen der Unterrichtsreihe}\\[0.3cm]
    {\Large\textit{\reihenthema}}\\[2cm]

    \begin{tabular}{rl}
        \textbf{Fach:} & \fach{} (\gerniveau) \\[0.3cm]
        \textbf{Klasse:} & Jahrgangsübergreifend \jahrgangsstufe{} (Spätbeginner) \\[0.3cm]
        \textbf{Lehrwerk:} & \lehrwerk, \lehrwerkseite \\[0.3cm]
        \textbf{Datum:} & Dienstag, \unterrichtsdatum \\[0.3cm]
        \textbf{Zeit:} & Doppelstunde (90 Min.) \\[0.3cm]
        \textbf{Raum:} & --- \\[0.3cm]
    \end{tabular}

    \vfill

    {\large Lehrkraft: \lehrkraft}\\[0.3cm]
    {\normalsize\textcolor{gray}{\ausbildungsstand}}\\[1cm]

\end{titlepage}

% ============================================================================
% INHALTSVERZEICHNIS
% ============================================================================
\tableofcontents
\newpage

% ============================================================================
% 1. BEDINGUNGSANALYSE
% ============================================================================
\section{Bedingungsanalyse}

\subsection{Lerngruppe}
\begin{tabular}{ll}
    Kursgröße: & \lerngruppengroesse{} SuS (\maennlich{} m / \weiblich{} w / \divers{} d) \\
    Jahrgangsstufe: & \jahrgangsstufe{} (jahrgangsübergreifend) \\
    Sprachniveau: & A2 (Gruppe A) / B1 (Gruppe B) \\
    Fremdsprache: & 3. Fremdsprache (Spätbeginner ab Kl. 10) \\
\end{tabular}

\subsection{Gruppeneinteilung}

\textbf{Gruppe A} (7 SuS, A2-Niveau) -- Lehrwerk S. 31--32:
\begin{itemize}
    \item \textbf{E.H.} $\bigstar\bigstar\bigstar\bigstar$ -- TOP-Schülerin, Französisch-Hintergrund, 15 Punkte $\rightarrow$ \textit{Peer-Tutor für schwächere SuS}
    \item \textbf{L.N.} $\bigstar\bigstar\bigstar$ -- Proaktiv, aber etwas bequem $\rightarrow$ \textit{Durch aktivierende Aufgaben in Bewegung halten}
    \item \textbf{C.P.} $\bigstar\bigstar$ -- Gute Aussprache, 12 Punkte, ruhig
    \item \textbf{V.S.} $\bigstar\bigstar$ -- Konstant, fleißig, ruhig
    \item \textbf{L.K.} $\bigstar\bigstar$ -- Kl. 11 Nachzüglerin, 10 Punkte, noch catch-up
    \item \textbf{H.P.} $\bigstar$ -- Nachzügler, 10 Punkte
    \item \textbf{A.A.} $\bigstar$ -- Fleißig, aber unsicher
    \item \textbf{A.S.} $\bigstar$ -- Einziger Junge, 9 Punkte
\end{itemize}

\textbf{Gruppe B} (3 SuS, B1-Niveau) -- Lehrwerk S. 104:
\begin{itemize}
    \item Kl. 12 + stärkere Kl. 11 SuS
    \item Arbeiten selbstständig an Wortbildungsübungen
\end{itemize}

\subsection{Besonderheiten}
\begin{itemize}
    \item \textbf{Heterogenität:} Große Leistungsunterschiede (9--15 Punkte)
    \item \textbf{Jahrgangsübergreifend:} Kl. 10, 11 und 12 gemeinsam
    \item \textbf{Sprachstand:} A2 vs. B1 erfordert zweigleisige Planung
    \item \textbf{Gruppendynamik:} 1 Junge unter 9 Mädchen $\rightarrow$ A.S. bewusst einbinden
    \item \textbf{Peer-Tutoring:} E.H. kann schwächere SuS unterstützen
\end{itemize}

% ============================================================================
% 2. SACHANALYSE
% ============================================================================
\section{Sachanalyse}

\subsection{Gruppe A: Freizeitorte und Aktivitäten (S. 31--32, Teil C)}

\subsubsection{Sprachliche Strukturen}
\begin{itemize}
    \item \textbf{Wortschatz:} Freizeitorte (la piscina, el cine, el teatro, el parque, la discoteca, el centro comercial, la biblioteca, el polideportivo)
    \item \textbf{Fragestruktur:} ¿Adónde vamos? / ¿Adónde vas?
    \item \textbf{Verb:} \textit{ir a} + Ort (Voy a la piscina / Vamos al cine)
    \item \textbf{Kontraktionen:} a + el = al (aber: a la, a los, a las)
\end{itemize}

\subsubsection{Aufgabenstruktur S. 31--32}
\begin{enumerate}
    \item \textbf{Aufg. 1:} Plakat-Vokabeln erschließen -- Bilder den Wörtern zuordnen
    \item \textbf{Aufg. 2:} Hörübung -- Radio-Spots verstehen (globales Hörverstehen)
    \item \textbf{Aufg. 3:} Leseverstehen -- Martas E-Mail über Freizeitaktivitäten
\end{enumerate}

\subsection{Gruppe B: Wortbildung (S. 104)}

\subsubsection{Sprachliche Strukturen}
\begin{itemize}
    \item \textbf{Suffixe:}
    \begin{itemize}
        \item -ción: comunicar $\rightarrow$ comunicación, informar $\rightarrow$ información
        \item -(i)dad: responsable $\rightarrow$ responsabilidad, capaz $\rightarrow$ capacidad
        \item -ar: grupo $\rightarrow$ agrupar, motivo $\rightarrow$ motivar
    \end{itemize}
    \item \textbf{Wortfeld:} Berufe, abstrakte Begriffe, Eigenschaften
\end{itemize}

\subsubsection{Aufgabenstruktur S. 104}
\begin{enumerate}
    \item \textbf{Aufg. 4a:} Substantive auf -ción bilden
    \item \textbf{Aufg. 4b:} Substantive auf -(i)dad bilden
    \item \textbf{Aufg. 4c:} Verben auf -ar bilden
    \item \textbf{Aufg. 5:} ¿Qué hace un buen jefe? -- freies Sprechen
\end{enumerate}

\subsection{Kontrastive Betrachtung (Deutsch $\leftrightarrow$ Spanisch)}
\begin{itemize}
    \item \textbf{Positive Transfers:} Internationalismen (cine $\sim$ Kino, teatro $\sim$ Theater), Suffixe ähnlich zu lat./frz. (-tion $\sim$ -ción)
    \item \textbf{Interferenzen:} \textit{al} als Kontraktion oft vergessen (fehlerhaft: *a el cine)
    \item \textbf{Falsche Freunde:} biblioteca $\neq$ Bibliothek (gemeint ist Bücherei, nicht Studentenbibliothek)
\end{itemize}

% ============================================================================
% 3. DIDAKTISCHE ANALYSE
% ============================================================================
\section{Didaktische Analyse}

\subsection{Stellung in der Sequenz}
Dies ist die \textbf{2. Doppelstunde} der Beobachtungssequenz (4 Doppelstunden gesamt).

\begin{tabular}{|c|l|l|}
\hline
\textbf{DS} & \textbf{Gruppe A (A2)} & \textbf{Gruppe B (B1)} \\
\hline
1 & Wiederholung, Diagnose & Selbstständige Arbeit \\
\hline
\textbf{2} & \textbf{S. 31--32: Freizeitorte} & \textbf{S. 104: Wortbildung} \\
\hline
3 & Vertiefung, Anwendung & Textproduktion \\
\hline
4 & Präsentation, Reflexion & Präsentation, Reflexion \\
\hline
\end{tabular}

\subsection{Kompetenzziele (Rahmenplan MV)}
\begin{itemize}
    \item \textbf{Kommunikativ:}
    \begin{itemize}
        \item Hörverstehen (A): Radio-Spots verstehen (globales HV)
        \item Leseverstehen (A): E-Mail über Freizeit verstehen
        \item Sprechen (B): Diskussion über Führungsqualitäten
    \end{itemize}
    \item \textbf{Sprachlich:}
    \begin{itemize}
        \item Wortschatz Freizeitorte erweitern (A)
        \item Wortbildungsregeln anwenden (B)
        \item Verb \textit{ir a} + Ort korrekt verwenden (A)
    \end{itemize}
    \item \textbf{Interkulturell:}
    \begin{itemize}
        \item Freizeitgestaltung spanischer Jugendlicher kennenlernen
    \end{itemize}
    \item \textbf{Methodisch:}
    \begin{itemize}
        \item Erschließungsstrategien: Wörter aus Kontext/Bildern ableiten
        \item Wortfamilienbildung als Lernstrategie (B)
    \end{itemize}
\end{itemize}

\subsection{Legitimation}
\textbf{Rahmenplan Spanisch MV (Sek I/II):}
\begin{itemize}
    \item Kompetenzbereich: Funktionale kommunikative Kompetenz
    \item Standards: Hör-/Hörsehverstehen A2, Leseverstehen A2, Sprechen B1
    \item Sprachliche Mittel: Wortschatz themengebunden erweitern
\end{itemize}

% ============================================================================
% 4. STUNDENZIELE
% ============================================================================
\section{Stundenziele}

\subsection{Grobziel}
Die SuS erweitern ihre \textbf{kommunikative Kompetenz} im Bereich Hören, Lesen und Sprechen, indem sie:
\begin{itemize}
    \item (Gruppe A) Freizeitorte benennen und über Freizeitaktivitäten sprechen
    \item (Gruppe B) Wortbildungsregeln anwenden und über Führungsqualitäten diskutieren
\end{itemize}

\subsection{Feinziele (operationalisiert)}

\textbf{Gruppe A:}
\begin{tabular}{|c|p{7.5cm}|c|c|}
\hline
\textbf{Nr.} & \textbf{Die SuS können...} & \textbf{AFB} & \textbf{SuS} \\
\hline
FZ1 & mindestens 6 Freizeitorte korrekt benennen und aussprechen & I & alle \\
\hline
FZ2 & Radio-Spots verstehen und die genannten Orte identifizieren & I & alle \\
\hline
FZ3 & die E-Mail von Marta lesen und Fragen zum Text beantworten & II & alle \\
\hline
FZ4 & mit \textit{ir a} + Ort Sätze über eigene Freizeitpläne bilden & II & C.P., V.S., E.H., L.N. \\
\hline
\end{tabular}

\vspace{1em}

\textbf{Gruppe B:}
\begin{tabular}{|c|p{7.5cm}|c|c|}
\hline
\textbf{Nr.} & \textbf{Die SuS können...} & \textbf{AFB} & \textbf{SuS} \\
\hline
FZ5 & aus Verben/Adjektiven Substantive mit -ción / -(i)dad bilden & II & alle \\
\hline
FZ6 & aus Substantiven Verben auf -ar ableiten & II & alle \\
\hline
FZ7 & die Eigenschaften eines guten Chefs beschreiben und begründen & III & alle \\
\hline
\end{tabular}

% ============================================================================
% 5. METHODISCHE ANALYSE
% ============================================================================
\section{Methodische Analyse}

\subsection{Differenzierungskonzept}
\begin{itemize}
    \item \textbf{Parallele Gruppenarbeit:} Gruppe A und B arbeiten an unterschiedlichen Themen
    \item \textbf{Peer-Tutoring:} E.H. unterstützt schwächere SuS in Gruppe A
    \item \textbf{Tempo-Differenzierung:} Schnellere SuS helfen langsameren
    \item \textbf{Scaffolding Gruppe A:} Wortschatzhilfen, Bildunterstützung
    \item \textbf{Selbstständigkeit Gruppe B:} Arbeiten weitgehend eigenständig
\end{itemize}

\subsection{Phasierung (Doppelstunde = 2 × 45 Min.)}
\begin{enumerate}
    \item \textbf{1. Stunde:}
    \begin{itemize}
        \item Gemeinsamer Einstieg: ¿Qué haces en tu tiempo libre?
        \item Gruppe A: Aufg. 1 (Plakat-Vokabeln), Aufg. 2 (Hörübung)
        \item Gruppe B: Aufg. 4a--4c (Wortbildungsübungen)
    \end{itemize}
    \item \textbf{2. Stunde:}
    \begin{itemize}
        \item Gruppe A: Aufg. 3 (Leseverstehen), Anwendung \textit{ir a}
        \item Gruppe B: Aufg. 5 (Sprechübung: Guter Chef)
        \item Gemeinsamer Abschluss: Ergebnissicherung
    \end{itemize}
\end{enumerate}

\subsection{Medien}
\begin{itemize}
    \item Tafel/Whiteboard
    \item Lehrwerk: A\_tope.com, S. 31--32 (A) / S. 104 (B)
    \item Audio: Hörtrack zu Aufg. 2 (Radio-Spots)
    \item Arbeitsblatt Wortbildung (für Gruppe B)
    \item Bildkarten Freizeitorte (optional für Gruppe A)
\end{itemize}

% ============================================================================
% 6. VERLAUFSPLANUNG
% ============================================================================
\section{Verlaufsplanung}

\textbf{1. Stunde (45 Min.)}

\begin{longtable}{|p{1cm}|p{1.5cm}|p{3.8cm}|p{3.2cm}|p{1.5cm}|p{2cm}|}
\hline
\textbf{Zeit} & \textbf{Phase} & \textbf{Lehrerhandeln} & \textbf{Schülerhandeln} & \textbf{SF} & \textbf{Medien} \\
\hline
\endhead

0--5 & Einstieg & Begrüßt, fragt: ¿Qué haces en tu tiempo libre? Sammelt Antworten an der Tafel & Nennen Freizeitaktivitäten (bekannter WS) & Plenum & Tafel \\
\hline

5--8 & Über\-leitung & Erklärt Gruppenaufteilung: Gruppe A S. 31, Gruppe B S. 104. Verteilt AB an Gruppe B & Notieren Seitenzahlen, B erhält AB & Plenum & Buch, AB \\
\hline

8--18 & Erarbei\-tung A & \textbf{Gruppe A:} Leitet Aufg. 1 an -- Plakat betrachten, Vokabeln erschließen. Hilft bei Aussprache. \newline \textit{Gruppe B arbeitet selbstständig an Aufg. 4a} & \textbf{A:} Ordnen Bilder den Wörtern zu, sprechen nach. \newline \textbf{B:} Bearbeiten 4a (Subst. auf -ción) & EA/PA & Buch S. 31, AB \\
\hline

18--28 & HV A & \textbf{Gruppe A:} Spielt Hörtext Aufg. 2 (Radio-Spots) 2x ab. Fragt: ¿Adónde van? \newline \textit{Gruppe B: Aufg. 4b} & \textbf{A:} Hören, notieren Orte, vergleichen mit Partner. \newline \textbf{B:} Bearbeiten 4b (Subst. auf -(i)dad) & EA $\rightarrow$ PA & Audio, Buch \\
\hline

28--35 & Sicherung A & \textbf{Gruppe A:} Bespricht Lösungen Aufg. 2. Klärt unbekannte Vokabeln. \newline \textit{Gruppe B: Aufg. 4c} & \textbf{A:} Vergleichen, korrigieren. \newline \textbf{B:} Bearbeiten 4c (Verben auf -ar) & Plenum A & Tafel \\
\hline

35--40 & Wort\-schatz A & Schreibt Vokabeln an die Tafel: la piscina, el cine, el teatro... Chorisches Sprechen & \textbf{A:} Sprechen nach, notieren ins Vokabelheft. \newline \textbf{B:} Kontrollieren 4a--c mit Lösungen & Plenum A & Tafel, Vokabelheft \\
\hline

40--45 & Pause / Über\-leitung & Kurze Bewegungspause. Ankündigung 2. Stunde: Gruppe A liest E-Mail, Gruppe B spricht über Chefs & Kurze Pause, Materialsichtung & --- & --- \\
\hline
\end{longtable}

\vspace{1em}

\textbf{2. Stunde (45 Min.)}

\begin{longtable}{|p{1cm}|p{1.5cm}|p{3.8cm}|p{3.2cm}|p{1.5cm}|p{2cm}|}
\hline
\textbf{Zeit} & \textbf{Phase} & \textbf{Lehrerhandeln} & \textbf{Schülerhandeln} & \textbf{SF} & \textbf{Medien} \\
\hline
\endhead

0--5 & Wieder\-holung & Schnelles Vokabelquiz Freizeitorte (Gruppe A). Fragt Gruppe B nach Ergebnissen 4a--c & \textbf{A:} Antworten auf Quiz. \newline \textbf{B:} Nennen 2--3 Beispiele & Plenum & --- \\
\hline

5--15 & LV A & \textbf{Gruppe A:} Leitet Aufg. 3 an -- Martas E-Mail lesen. E.H. hilft schwächeren SuS. \newline \textit{Gruppe B: Bereiten Aufg. 5 vor} & \textbf{A:} Lesen still, markieren Schlüsselwörter. \newline \textbf{B:} Notieren Stichpunkte für Diskussion & EA & Buch S. 32 \\
\hline

15--22 & Sicherung LV & \textbf{Gruppe A:} Fragt Verständnisfragen zur E-Mail: ¿Adónde va Marta? ¿Con quién? ¿Cuándo? & \textbf{A:} Beantworten mündlich, vergleichen & PA $\rightarrow$ Plenum A & Buch \\
\hline

22--32 & Anwen\-dung & \textbf{Gruppe A:} Partnerarbeit -- Mini-Dialoge mit \textit{ir a}: ¿Adónde vamos el sábado? Vamos al/a la... \newline \textbf{Gruppe B:} Aufg. 5 in Kleingruppe & \textbf{A:} Üben Dialoge. E.H. als Partnerin für schwächere SuS. \newline \textbf{B:} Diskutieren: ¿Qué hace un buen jefe? & PA & --- \\
\hline

32--38 & Präsen\-tation & 2--3 Paare aus Gruppe A präsentieren ihren Dialog. Gruppe B fasst Diskussionsergebnis zusammen & \textbf{A:} Präsentieren Dialoge. \newline \textbf{B:} Ein/e SuS fasst zusammen & Plenum & --- \\
\hline

38--42 & Sicherung & Fasst zusammen: Neue Vokabeln, Struktur \textit{ir a}, Wortbildungsregeln. Schreibt Merkpunkte an Tafel & Notieren wichtige Punkte & Plenum & Tafel \\
\hline

42--45 & Abschluss & Gibt HA auf: Gruppe A: Vokabeln lernen. Gruppe B: 5 Sätze mit neuen Wörtern schreiben & Notieren HA & Plenum & --- \\
\hline
\end{longtable}

\textbf{Legende:} EA = Einzelarbeit, PA = Partnerarbeit, GA = Gruppenarbeit, SF = Sozialform, LV = Leseverstehen, HV = Hörverstehen

% ============================================================================
% 7. ERWARTETE SCHWIERIGKEITEN
% ============================================================================
\section{Erwartete Schwierigkeiten}

\begin{tabular}{|p{5cm}|p{8cm}|}
\hline
\textbf{Schwierigkeit} & \textbf{Reaktion} \\
\hline
Kontraktion \textit{al} vergessen (*a el cine) & Explizite Kontrastierung: a + el = al, a + la bleibt a la. Regel an Tafel. Chorisches Sprechen beider Formen \\
\hline
Hörverstehen zu schnell & Audio 2x abspielen. Vor 2. Hören: Vorentlastung der Schlüsselwörter. Ggf. pausieren \\
\hline
Gruppe A langsamer als geplant & Puffer: Anwendungsphase (Dialoge) kann gekürzt werden. E.H. als Peer-Tutor beschleunigt \\
\hline
Gruppe B zu schnell fertig & Zusatzaufgabe: Weitere 5 Beispiele pro Wortbildungsregel finden. Oder: Unterstützung Gruppe A \\
\hline
A.S. isoliert als einziger Junge & Bewusst in gemischte Partnerarbeit einteilen. Bei Dialogen mit L.N. oder C.P. paaren \\
\hline
H.P. und A.A. unsicher bei Aufg. 3 & E.H. unterstützt als Peer-Tutor. Vereinfachte Fragen vorbereiten \\
\hline
L.N. unterfordert / bequem & Aktivierende Rolle: Soll Dialog vorstellen, Vokabeln an Tafel schreiben \\
\hline
\end{tabular}

% ============================================================================
% 8. HAUSAUFGABE
% ============================================================================
\section{Hausaufgabe}

\textbf{Gruppe A (A2):}
\begin{itemize}
    \item Vokabeln S. 31--32 lernen (Freizeitorte)
    \item Optional: Eigene E-Mail über Freizeitpläne schreiben (3--5 Sätze mit \textit{ir a})
\end{itemize}

\textbf{Gruppe B (B1):}
\begin{itemize}
    \item 5 Sätze mit neu gebildeten Wörtern aus Aufg. 4 schreiben
    \item Vokabeln Wortbildung wiederholen
\end{itemize}

% ============================================================================
% 9. ANLAGEN
% ============================================================================
\section{Anlagen}

\begin{enumerate}
    \item Lehrwerk A\_tope.com, S. 31--32 (Gruppe A: Freizeitorte)
    \item Lehrwerk A\_tope.com, S. 104 (Gruppe B: Wortbildung)
    \item Audio-Track Aufg. 2 (Radio-Spots)
    \item Arbeitsblatt Wortbildung (für Gruppe B) -- ggf. noch zu erstellen
    \item Tafelbild: Freizeitorte mit Artikeln, Regel \textit{a + el = al}
\end{enumerate}

\subsection{Tafelbild (Entwurf)}

\begin{verbatim}
+------------------------------------------------------------------+
|  ¿Adónde vamos? - Freizeitorte                                   |
+------------------------------------------------------------------+
|                                                                  |
|  la piscina         el cine           el teatro                  |
|  (Schwimmbad)       (Kino)            (Theater)                  |
|                                                                  |
|  el parque          la discoteca      el centro comercial        |
|  (Park)             (Disco)           (Einkaufszentrum)          |
|                                                                  |
|  la biblioteca      el polideportivo                             |
|  (Bibliothek)       (Sportzentrum)                               |
|                                                                  |
+------------------------------------------------------------------+
|  REGLA:  a + el = al      |  Voy AL cine.                        |
|          a + la = a la    |  Voy A LA piscina.                   |
+------------------------------------------------------------------+
\end{verbatim}

% ============================================================================
% 10. DIDAKTIK-SCORE
% ============================================================================
\newpage
\section{Didaktik-Score}

\begin{scorebox}
\textbf{Bewertung der didaktischen Qualität nach 10 Dimensionen}\\
\textbf{Ziel-Score: $\geq$ 9.0 (Premium-Qualität)}

\vspace{1em}
\begin{tabular}{|c|l|c|c|p{4.5cm}|}
\hline
\textbf{\#} & \textbf{Dimension} & \textbf{Gew.} & \textbf{Score} & \textbf{Notiz} \\
\hline
1 & Differenzierung & 15\% & 9/10 & Zwei Niveaugruppen mit unterschiedlichen Inhalten; Peer-Tutoring \\
\hline
2 & Taxonomie (AFB) & 10\% & 9/10 & AFB I (Vokabeln), II (Anwendung), III (Diskussion B) \\
\hline
3 & Lernziel-Alignment & 10\% & 10/10 & Ziele direkt aus Rahmenplan, Lehrwerk-kongruent \\
\hline
4 & Aktivierung & 10\% & 9/10 & HV, LV, Dialoge, Diskussion -- verschiedene Fertigkeiten \\
\hline
5 & Scaffolding & 10\% & 9/10 & Bildunterstützung, Peer-Tutor, gestufte Hilfen \\
\hline
6 & Feedback-Qualität & 10\% & 8/10 & Mündliches Feedback; Lösungsvergleich; wenig schriftlich \\
\hline
7 & Sprachliche Klarheit & 10\% & 10/10 & Klare Aufgabenstellungen im Lehrwerk \\
\hline
8 & Motivation \& Relevanz & 10\% & 9/10 & Alltagsthema Freizeit; Berufsthema für B1 relevant \\
\hline
9 & Inklusion (UDL) & 10\% & 8/10 & Visuell (Bilder), auditiv (HV), kinästhetisch könnte mehr sein \\
\hline
10 & Praktikabilität & 5\% & 9/10 & 90 Min. realistisch; Material vorhanden \\
\hline
\multicolumn{3}{|r|}{\textbf{GESAMT:}} & \textbf{9.0/10} & \textcolor{successgreen}{\textbf{FREIGABE}} \\
\hline
\end{tabular}

\vspace{1em}
\textbf{Bewertungsschwellen:}
\begin{itemize}
    \item[\checkmark] \textcolor{successgreen}{$\geq$ 9.0: \textbf{FREIGABE} (Premium-Qualität)}
    \item[$\Box$] \textcolor{warningorange}{8.0--8.9: Iteration empfohlen}
    \item[$\Box$] 6.0--7.9: Überarbeitung erforderlich
    \item[$\Box$] $<$ 6.0: Grundlegende Neukonzeption
\end{itemize}

\vspace{1em}
\textbf{Gewichtete Berechnung:}
$(9 \times 0.15) + (9 \times 0.10) + (10 \times 0.10) + (9 \times 0.10) + (9 \times 0.10) + (8 \times 0.10) + (10 \times 0.10) + (9 \times 0.10) + (8 \times 0.10) + (9 \times 0.05)$
$= 1.35 + 0.90 + 1.00 + 0.90 + 0.90 + 0.80 + 1.00 + 0.90 + 0.80 + 0.45 = \textbf{9.00}$

\end{scorebox}

\end{document}
